% =====================================================================================
% ANNEXE AJOUTEE LE 12 AOUT 2026, a la demande de Caroline Hillairet, qui a salue
% l'initiative de pallier le manque de donnees par un recueil d'avis d'experts et demande de
% mettre la methode de Cooke et le questionnaire en annexe.
%
% LE SENS DE L'ANNEXE A CHANGE AVEC LA DECISION DE NE PAS LANCER L'ELICITATION. Elle ne
% presente donc pas un protocole A VENIR mais un protocole PREPARE ET NON EXECUTE, avec le
% chiffrage de la raison. C'est un actif devant un jury qui demanderait pourquoi un modele dont
% la direction n'est pas identifiable ne recourt pas au jugement d'expert : la reponse existe,
% elle est chiffree, et le materiel est pret.
%
% ATTENTION AU PIEGE DU HARNAIS, ET IL A ETE RENCONTRE ICI. Une declaration
% HARNAIS-HORS-SCRIPT posee AVANT la premiere \section couvre tout le fichier : la premiere
% version de cette annexe declarait ainsi ses onze nombres d'un coup, dont les huit qui sont
% parfaitement verifiables. La declaration est donc descendue dans la seule section qui la
% justifie, et les autres citent leurs scripts.
% =====================================================================================

\chapter{Le protocole d'élicitation, préparé et non exécuté}
\label{ch:elicitation}

% SOURCES-SCRIPTS: 31 34 59

La direction de la matrice de contagion n'est pas identifiable sur la donnée disponible
(chapitre~\ref{ch:identifiabilite}). Le recours au jugement d'expert formalisé est la réponse
classique à ce genre d'impasse, et un protocole complet a été construit pour cela. Il n'a pas
été exécuté, et cette annexe existe pour que ce choix soit lisible : le protocole y est décrit,
la raison de ne pas le lancer y est chiffrée, et ce que l'on perd y est dit.

\section{La méthode de Cooke, en quatre temps}

% HARNAIS-HORS-SCRIPT: les quantiles demandes a l'expert sont une propriete de CONCEPTION du
% questionnaire, non une grandeur calculee : aucun script ne les imprime et aucun ne le devrait.

Le \emph{modèle classique} de Cooke ne traite pas les experts comme interchangeables : il les
\emph{note}, sur des questions dont l'analyste connaît la réponse, puis pondère leurs avis par
cette note. Il se déroule en quatre temps.

\begin{enumerate}[itemsep=3pt,topsep=3pt]
\item \textbf{Élicitation en quantiles.} Chaque expert ne donne pas un point mais trois
      quantiles, à $5$, $50$ et $95\,\%$, ce qui rend son incertitude déclarée mesurable.
\item \textbf{Calibration.} Sur les \emph{questions-graines}, dont la réponse est connue de
      l'analyste et masquée à l'expert, on compare la fréquence empirique des réalisations dans
      les intervalles déclarés à la fréquence théorique. L'écart se lit comme une statistique
      d'information de Kullback-Leibler, et son niveau de signification devient le
      score de calibration.
\item \textbf{Information.} Un expert bien calibré mais qui déclare des intervalles très larges
      n'apporte rien. Le score d'information mesure la concentration de ses
      distributions relativement à une loi de référence.
\item \textbf{Poids et agrégation.} Le poids d'un expert est le produit des deux scores, et il
      est nul sous un seuil de calibration. L'avis agrégé est la combinaison linéaire
      pondérée des distributions individuelles.
\end{enumerate}

C'est le quatrième point qui fait la force et la fragilité de la méthode : un expert mal
calibré est écarté, ce qui protège des dires péremptoires, mais réduit d'autant la taille
effective du panel.

\section{Le protocole construit pour la matrice de contagion}

Le questionnaire préparé comporte deux parties. La partie d'étalonnage pose des
questions à réponse connue, tirées de sources publiques vérifiables, sur des grandeurs du même
domaine que les cibles. La partie cible interroge les liens dirigés entre piliers, un
par un, sous la forme d'une probabilité conditionnelle : sachant que le pilier source est
défaillant, quelle est la probabilité que le pilier cible le devienne.

Le dispositif complet demande seize questions par expert, dont dix questions-graines. Ce ratio
n'est pas un excès de prudence : c'est le minimum qui permette de distinguer un expert calibré
d'un expert chanceux, et le réduire viderait la méthode de ce qui la distingue d'une moyenne
d'opinions. Le protocole, le questionnaire et le gabarit de réponses sont versionnés dans le
dépôt du mémoire, et deux scripts implémentent l'agrégation, l'un sur réponses fictives pour
valider la chaîne de traitement, l'autre prêt à recevoir des réponses réelles.

\paragraph{Une dixième graine a été retirée, et le dire fait partie du protocole.} La première
version de l'étalonnage portait une question sur l'indice de queue de sévérité $\xi$, de valeur
vraie annoncée $0{,}90$. Deux défauts s'y cumulaient. D'abord $0{,}90$ n'est pas la calibration
mais la valeur \emph{posée} du scénario de pire cas, la calibration figée donnant
$\hat\xi = 0{,}5954$ : or dans la méthode de Cooke une graine fausse n'ajoute pas du bruit, elle
\emph{inverse} les poids, en pénalisant le répondant qui vise juste. Ensuite, et c'est ce qui
justifie le retrait plutôt que la correction, $\xi$ est une \emph{estimation} d'intervalle
$[0{,}30\,;0{,}83]$, plus large que la valeur elle-même : une graine doit être une grandeur
réalisée, pas un paramètre estimé. Les dix graines de l'instrument révisé sont toutes reproductibles par une
sortie de script versionnée, ce qui est la condition pour qu'une notation d'expert soit
opposable.

\vspace{2pt}
% SOURCES-SCRIPTS: 26 47 63

%--------------------------------------------------------------
\section{Le questionnaire, reproduit}
\label{sec:questionnaire-reproduit}
%--------------------------------------------------------------
% LES VALEURS VRAIES DES GRAINES NE SONT PAS PUBLIEES ICI, ET C'EST DELIBERE : les imprimer
% BRULERAIT l'instrument, un expert qui a lu le memoire ne pouvant plus etre note dessus. Le
% script 105 les calcule, leur applique quatre controles et les imprime dans sa sortie de
% travail ; si l'elicitation etait lancee apres publication, les graines se retirent du meme jeu
% de donnees par ce meme script. Ce qui est cite ci-dessous est donc le seul nombre que
% l'arbitrage de conception exige de montrer, la valeur du candidat ecarte.
% SOURCES-SCRIPTS: 105

Décrire un instrument ne suffit pas à établir qu'il existe. Il est donc reproduit ici dans la
forme sous laquelle il aurait été soumis, de façon qu'un lecteur puisse juger sur pièce ce que
la section précédente affirme, et le réutiliser si la décision était rouverte. Le formulaire
complet, avec sa consigne, son exemple traité, sa mention d'usage et de conservation et son
accord signé du répondant, est versionné à part dans le dépôt du mémoire ; la présente annexe en
reprend les deux tables et l'essentiel de la consigne.

\begin{attention}
\textbf{L'instrument reproduit ici est la version révisée, et il diffère de l'exemplaire de
travail.} Celui-ci portait une question d'étalonnage sur l'indice de queue de sévérité, retirée
pour les deux motifs exposés plus haut. Elle n'est pas corrigée mais \emph{remplacée}, et le
remplacement a lui-même été arbitré : un indice de Gini de la concentration avait d'abord été
retenu, puis écarté parce qu'il fait doublon avec la question qui compte les incidents portant la
moitié du volume, et parce qu'une valeur vraie de $0{,}97$ se devine sans rien connaître du
dossier. La graine retenue porte sur la durée de confinement, qui est une performance de
résilience et donc du même domaine que les cibles. \textbf{L'exemplaire de travail ne doit pas
être soumis en l'état.}
\end{attention}

\paragraph{La règle que s'impose l'instrument, et elle est vérifiée.} Une graine doit être une
\emph{grandeur réalisée}, reproductible par une sortie de script versionnée, et jamais un
paramètre estimé. Les dix graines retenues sont donc des comptages, des parts, des délais ou des
quantiles observés : aucune ne demande d'ajuster un modèle. Un script les calcule et leur applique
quatre contrôles, portant sur ce caractère réalisé, sur la non-dégénérescence de la valeur, sur le
respect des bornes naturelles, et sur le pouvoir de séparation entre experts. Une seule graine est
signalée comme potentiellement devinable, celle qui porte sur la part du piratage, et elle est
conservée avec son motif : sa valeur vraie est nettement au-dessus de la réponse spontanée, si
bien qu'elle sépare malgré son caractère extrême.

\paragraph{Aucune graine ne porte une grandeur de capital, et c'est délibéré.} Un praticien de la
résilience observe des incidents, des délais et des concentrations ; il n'observe pas un quantile
de charge annuelle. Noter un expert sur une telle grandeur mesurerait sa familiarité avec le
présent mémoire, non son jugement sur le risque.

\subsection*{Consigne remise à l'expert}

\begin{cle}
\textbf{Comment répondre.} Pour chaque quantité, donnez trois nombres : une valeur \emph{basse}
($5\,\%$), une valeur \emph{médiane} ($50\,\%$), une valeur \emph{haute} ($95\,\%$), telles que
vous estimez à $90\,\%$ de chances que la vraie valeur soit comprise entre votre $5\,\%$ et votre
$95\,\%$. Il n'y a pas de réponse attendue : donnez votre meilleure estimation \emph{et} votre
incertitude. Répondez même incertain, un intervalle large étant une réponse honnête et
parfaitement valable. Ne cherchez pas les vraies valeurs : l'objet est de mesurer votre jugement.
\end{cle}

\subsection*{Partie A, les dix questions d'étalonnage}

Des quantités observées dans les bases du mémoire, dont la valeur vraie est connue de l'analyste
et n'est pas affichée au répondant. Elles servent à pondérer les experts, non à les piéger. Les
trois colonnes de droite sont vides dans les deux tables qui suivent : c'est l'espace de réponse
du formulaire, non une donnée manquante.

{\small
\renewcommand{\arraystretch}{1.22}
\begin{longtable}{@{}l p{8.0cm} c c c@{}}
\toprule
& \textbf{Quantité, sur les incidents notifiés de la chronologie} & $\mathbf{5\,\%}$ &
$\mathbf{50\,\%}$ & $\mathbf{95\,\%}$ \\
\midrule
\endfirsthead
\toprule
& \textbf{Quantité (suite)} & $\mathbf{5\,\%}$ & $\mathbf{50\,\%}$ & $\mathbf{95\,\%}$ \\
\midrule
\endhead
G1 & Délai médian entre la survenance d'un incident et sa déclaration, en jours & & & \\
G2 & Part des incidents déclarés plus de quatre-vingt-dix jours après leur survenance,
     en \% & & & \\
G3 & Délai médian de déclaration d'un incident de type piratage, en jours & & & \\
G4 & Durée médiane entre le début et la fin d'exposition d'un incident, en jours & & & \\
G5 & Nombre médian d'enregistrements touchés par incident & & & \\
G6 & Nombre d'incidents qui portent à eux seuls la moitié du volume touché
     sur sept exercices & & & \\
G7 & Part des incidents dont l'exposition a duré plus de trente jours, en \% & & & \\
G8 & Part des incidents de type piratage parmi ceux dont le vecteur est renseigné,
     en \% & & & \\
G9 & Part des lignes de la chronologie qui sont des doublons de notification d'un même
     incident, en \% & & & \\
G10 & Rapport du quatre-vingt-dix-neuvième centile à la médiane des pertes cyber du
      secteur financier & & & \\
\bottomrule
\end{longtable}}

\vspace{2pt}
\noindent\textbf{Ce que chaque graine cherche à mesurer.} G1 à G3 portent sur le \emph{délai de
détection}, qui est le canal que la conformité déplace le plus directement. G4 et G7 portent sur
le \emph{confinement}. G5, G6 et G10 portent sur la \emph{concentration} de la perte, c'est-à-dire
sur l'intuition de queue lourde. G8 porte sur la structure des vecteurs et G9 sur la qualité de la
donnée déclarative, dont un praticien de la notification a une expérience directe.

\subsection*{Partie B, les six liens dirigés entre piliers}

Force du lien dirigé de $j$ vers $k$, entre $0$ et $1$ : $0$ signifie qu'une défaillance du
pilier $j$ n'entraîne jamais celle de $k$, $1$ qu'elle l'entraîne quasi certainement dans la
foulée. Les trois mêmes quantiles sont demandés.

\begin{attention}
\textbf{Le sens des lettres n'est pas une coquette\-rie de notation.} La matrice du modèle
s'écrit $W_{jk}$ avec $j$ la \emph{source} et $k$ la \emph{cible} : la ligne émet, la colonne
reçoit, et deux scripts vérifient cette convention avant tout calcul. Un formulaire qui
demanderait la force du lien \emph{de $k$ vers $j$} ferait donc saisir la matrice
\emph{transposée}. Or c'est précisément entre une matrice et sa transposée que le
chapitre~\ref{ch:identifiabilite} montre que la donnée ne tranche pas, et le
chapitre~\ref{chap:cascade} que le capital et la décision diffèrent. Une inversion de lettres
dans l'instrument produirait donc exactement l'erreur que tout le mémoire s'emploie à ne pas
commettre, et elle ne serait rattrapée par aucun contrôle numérique.
\end{attention}

{\small
\renewcommand{\arraystretch}{1.25}
\begin{longtable}{@{}l p{8.2cm} c c c@{}}
\toprule
& \textbf{Lien dirigé $j \rightarrow k$, $j$ émet} & $\mathbf{5\,\%}$ & $\mathbf{50\,\%}$ &
$\mathbf{95\,\%}$ \\
\midrule
\endfirsthead
\toprule
& \textbf{Lien dirigé (suite)} & $\mathbf{5\,\%}$ & $\mathbf{50\,\%}$ & $\mathbf{95\,\%}$ \\
\midrule
\endhead
B1 & P1 $\rightarrow$ P2 : une défaillance de gouvernance entraîne des incidents & & & \\
B2 & P2 $\rightarrow$ P1 : des incidents entraînent une défaillance de gouvernance & & & \\
B3 & P4 $\rightarrow$ P2 : une défaillance d'un tiers entraîne des incidents & & & \\
B4 & P1 $\rightarrow$ P4 : une défaillance de gouvernance entraîne une défaillance côté
     tiers & & & \\
B5 & P3 $\rightarrow$ P1 : une défaillance des tests de résilience entraîne une défaillance
     de gouvernance & & & \\
B6 & P5 $\rightarrow$ P2 : une défaillance du partage d'information entraîne des
     incidents & & & \\
\bottomrule
\end{longtable}}

\paragraph{Le couple B1 et B2 est l'objet même du mémoire.} Les deux questions portent sur le
\emph{même} couple de piliers, dans les deux sens. Une donnée de co-occurrence ne les distingue
pas ; un expert, lui, peut les séparer. C'est précisément la quantité que le
chapitre~\ref{ch:identifiabilite} montre non identifiable, et c'est la raison pour laquelle ce
questionnaire aurait eu une valeur que nulle autre source du dossier ne porte.

\paragraph{Ce qui est fait des réponses.} La partie~A mesure la calibration du répondant, c'est-à-dire
la justesse de ses intervalles, et son information, c'est-à-dire leur finesse ; ces deux
grandeurs fixent son poids dans l'agrégation. La partie~B, pondérée par ce poids et croisée avec
les réponses des autres experts, produirait chaque lien de la matrice sous forme d'\emph{intervalle}
et non de point. Le traitement est automatique et aucune donnée d'incident n'est requise, seulement
le jugement.

\section{Pourquoi il n'a pas été exécuté}

La raison est arithmétique, et elle porte sur la \emph{taille effective} du panel.

\paragraph{Un panel de deux experts n'a aucun pouvoir de calibration.} Avec le calendrier du
mémoire et un taux de réponse réaliste sur un questionnaire de seize questions quantitatives, le
nombre de réponses exploitables, c'est-à-dire complètes et passant le seuil de calibration,
serait de l'ordre de une à deux. Or à deux experts la pondération de Cooke dégénère : soit un
seul survit au seuil et l'on obtient un avis unique déguisé en agrégation, soit les deux
survivent et le poids relatif est estimé sur dix graines, ce qui ne le distingue pas d'une
moyenne simple. Dans les deux cas la méthode ne fait plus ce pour quoi elle est choisie.

\paragraph{Et une élicitation faible ne serait pas neutre, elle serait trompeuse.} C'est
l'argument décisif, et il est chiffré. Le chapitre~\ref{chap:partielle} mesure l'effet d'un
panel biaisé sur le capital : il déplace le \textsc{scr} de $777$~M\euro{} sans
qu'aucun signal ne l'indique au lecteur, là où l'approche par bornes reste \emph{invariante} à
ce biais. Un jugement d'expert faiblement calibré introduirait donc dans le résultat une
incertitude \emph{non déclarable}, alors que la non-identifiabilité, elle, se déclare et se
borne. En termes de qualité de la conclusion, l'échange serait défavorable : on remplacerait
une ignorance mesurée par une erreur invisible.
% SOURCES-SCRIPTS: 31

\paragraph{Une asymétrie à assumer, car un relecteur la trouvera.} Cet argument condamne le
jugement d'expert non vérifiable, or ce mémoire \emph{utilise} du jugement d'expert ailleurs :
la direction de la contagion s'appuie sur un corpus de dix rapports codés, et par un seul codeur
à ce jour. Refuser l'élicitation au nom du biais d'expert tout en s'appuyant sur un codage
d'expert demande donc une justification, et non un silence.

Elle tient en trois différences, et c'est leur cumul qui compte. Le codage est
\emph{auditable} : le protocole est écrit, les dix récits sont publics, et n'importe qui peut
recoder le corpus et confronter ses réponses aux nôtres, ce qu'une élicitation quantitative en
séance ne permet pas. Son biais principal est mesuré plutôt que redouté : le biais de
narration fait l'objet d'un test contre une loi nulle exacte, dont le résultat est publié
\emph{et négatif} (\S\ref{sec:biais-narration}). Et son produit est un ordre partiel,
un objet qualitatif que la suite du raisonnement borne, là où une élicitation produirait des
coefficients numériques entrant directement dans le capital. Un dire d'expert qui alimente une
borne et un dire d'expert qui alimente un nombre n'engagent pas le même risque.

Reste que la comparaison n'est pas entièrement favorable : le codage demeure à juge
unique, et cela se lit comme une limite déclarée et non comme une étape en attente. Ce
qui la distingue du défaut reproché à l'élicitation est qu'elle est \emph{bornée} : la
dégradation complète de la source contestée coûte moins que l'ignorance directionnelle déjà
publiée (\S\ref{sec:biais-narration}), là où un panel faiblement calibré introduirait une
incertitude non déclarable. Un second codage en aveugle renforcerait l'énoncé, et le dispositif
reste préparé à cette fin.

\section{Ce qui le remplace}

Trois dispositifs, tous en place et tous sans prior d'expert.

\begin{description}[itemsep=3pt,topsep=3pt]
\item[Le corpus documentaire.] Dix rapports post-mortem officiels sont codés selon un protocole
      écrit, et le sens des transitions y est testé contre une loi nulle exacte de permutation :
      $z = +5{,}12$ sur le corpus étendu. La direction n'est donc pas seulement posée, elle est
      \emph{corroborée} par une source indépendante du classeur qualitatif.
\item[L'échelle de repli sur la direction.] Plutôt que de connaître la direction, on peut
      supposer qu'elle est de rang faible. Le premier cran de cette échelle capte $74{,}1\,\%$
      de la variance de la direction posée et resserre le capital dans
      $[7\,761\,;\,8\,316]$~M\euro, soit une largeur de $555$~M\euro{} contre $1\,839$ sous
      ignorance totale. C'est une hypothèse structurelle, déclarée comme telle, et bien plus
      modeste qu'un dire d'expert.
\item[Les ancrages réglementaires.] Les niveaux d'exigence des textes et les publications des
      autorités fournissent des repères externes sur les états de conformité, sans passer par un
      panel.
\end{description}

% SOURCES-SCRIPTS: 30 34 59

\section{La relecture de praticien et son résultat}
\label{sec:relecture-praticien}

\paragraph{Pourquoi ce dispositif est nécessaire, et il l'est pour une raison qui n'a rien de
conjoncturel.} La direction de $W$ n'est pas identifiable sur les données disponibles,
et le chapitre~\ref{ch:identifiabilite} le démontre plutôt que de le déplorer : les statistiques
que les sources permettent de construire sont aveugles à la partie antisymétrique, donc une
matrice et sa transposée y sont indiscernables. Il en résulte que $W$ n'est pas calibrée au sens où le sont la fréquence et
la sévérité. Sa structure est \emph{déduite} d'un corpus de rapports post-mortem, source qui
porte son propre confondant, le biais de narration d'une enquête officielle, borné au
\S\ref{sec:biais-narration} mais non levé.

Sur une structure ainsi posée, il ne reste qu'un seul contrôle externe possible : le
jugement de praticiens qui ont vu des défaillances réelles, et qui n'ont lu ni le corpus ni le
mémoire. Ce dispositif est donc nécessaire et non facultatif. Mais il faut voir ce
qu'il peut et ce qu'il ne peut pas : un praticien peut corroborer ou contredire une
structure, il ne peut pas la calibrer, faute de quoi on retomberait sur l'élicitation
que ce chapitre vient d'écarter, avec son déplacement de capital non déclarable. C'est
exactement la différence entre une relecture et une élicitation, et elle commande la forme de
l'instrument.

Le dispositif est donc de nature différente des trois précédents : non pas recueillir des
probabilités, mais soumettre à des praticiens du secteur les affirmations \emph{qualitatives}
sur lesquelles le modèle repose, en leur demandant de les contredire. L'instrument est
un questionnaire de sept phrases, sans lecture préalable du mémoire, dont deux portent
explicitement sur l'ordre P1 vers P2. Le choix de demander des désaccords plutôt que des
accords est délibéré : un accord ne discrimine rien, alors qu'un désaccord argumenté désigne
l'endroit exact où le modèle s'écarte de ce que le terrain observe.

\paragraph{La première réponse ne contredit rien.} Une consultante de Nexialog, exerçant en
gestion des risques et du contrôle depuis une dizaine d'années, n'a contredit aucune
des sept phrases : deux accords pleins, cinq accords assortis d'une nuance. Un instrument
construit pour recueillir des désaccords et qui n'en recueille aucun renseigne d'abord sur
lui-même, et cette information vaut d'être écrite plutôt que passée sous silence.

\paragraph{Ce que la réponse corrobore, et ce n'est pas la direction.} Deux affirmations
reçoivent un accord franc et argumenté, et toutes deux portent sur la \emph{faisabilité} des
recommandations plutôt que sur la structure du modèle. Sur les trois indicateurs pilotables,
la répondante les juge cohérents et produisibles, et suggère de les compléter par une fréquence
d'occurrence. Sur le registre d'incidents mieux spécifié, qui est le prolongement le plus utile
identifié par ce mémoire, elle va plus loin que l'accord : elle décrit un usage professionnel
établi, un tableau de bord des incidents par prestataire servant, dans une analyse de risque
conduite au titre de DORA, à évaluer la fréquence du risque de défaillance système et le taux
d'indisponibilité sur douze mois. La recommandation la moins mathématique du mémoire est donc
aussi celle que le terrain confirme le plus directement.

\paragraph{Une qualification substantielle, qu'il faut déclarer sans la traiter.} La nuance la
plus utile porte sur l'arc de la gouvernance vers la gestion des prestataires tiers. La
répondante distingue l'occurrence et la maîtrise : une gouvernance solide ou
fragile ne change pas, selon elle, le fait qu'un prestataire critique tombe, mais elle change
la capacité de l'entité à en limiter les conséquences. Traduite dans les termes du modèle,
cette distinction déplacerait cet arc de la matrice de propagation vers le canal de détection,
ce qui n'est pas la même chose et ne donne pas le même chiffre. La réserve est donc
déclarée et non traitée : la traiter reviendrait à recalibrer, ce que le gel interdit,
et une réponse unique ne justifie pas de rouvrir la calibration.

\paragraph{Et un argument d'antériorité, le seul de la réponse qui soutienne vraiment un
ordre.} Sur la phrase qui fait de la gestion des incidents un puits, la répondante avance un
argument temporel plutôt qu'une intuition : les tests sont conduits en amont de la mise en
service, donc un incident postérieur ne peut pas dégrader les tests. C'est de la précédence, et
c'est précisément ce que le corpus documentaire code. Elle ajoute une critique de la définition
du pilier lui-même, en séparant l'incident de sa gestion, deux objets que la nomenclature
réglementaire réunit.

\paragraph{Ce que l'instrument devient.} Deux défauts du questionnaire ont été identifiés par
cette première réponse, tous deux dans les phrases et non chez la répondante. La phrase sur la
gouvernance source pure \emph{empaquetait} deux affirmations, et la répondante a nuancé la
première sans se prononcer sur la seconde, qui est celle qui compte. Et la phrase opposant la
contagion à la cause commune ne discriminait pas ce qu'elle prétendait discriminer, pour la
raison exposée au chapitre~\ref{ch:identifiabilite}. Trois questions à choix forcé ont été
ajoutées,
qui séparent les deux affirmations empaquetées, demandent au répondant de choisir entre
contagion et cause commune, et font trancher la distinction occurrence contre maîtrise. La
modalité \emph{je ne peux pas trancher} y est offerte explicitement, et c'est la plus
intéressante : un praticien qui la choisit corrobore la frontière d'identifiabilité sur le
terrain, ce qui vaut davantage qu'un accord sur la matrice.

\paragraph{La seconde réponse contredit deux phrases, et c'est ce que l'instrument cherchait.}
Un consultant de Nexialog, expert de la résilience et de la continuité d'activité dans les
services financiers depuis sept ans, a répondu à la même version du questionnaire, sans les
questions à choix forcé ajoutées depuis. Il ne se prononce pas sur la première phrase tout en la
commentant, donne quatre accords assortis d'une nuance, et exprime deux désaccords argumentés,
sur les deux phrases que le questionnaire désignait comme les plus importantes.

\paragraph{La cause commune, avancée sans y être invité.} Sur la phrase qui oppose la contagion à
la cause extérieure commune, il reconnaît des chaînes de causalité, une gouvernance peu
structurée conduisant à des tests insuffisants puis à des incidents, mais ajoute que plusieurs
domaines peuvent céder ensemble parce qu'ils partagent une cause : un manque de ressources, une
transformation mal pilotée, une architecture vieillissante. Il conclut qu'il parlerait de
dépendances entre domaines plutôt que de causalité systématique. Là où la première réponse
illustrait la frontière d'identifiabilité sans la nommer, celle-ci la formule : c'est la
corroboration la plus nette que le dispositif ait produite, et elle porte sur la frontière, non
sur la matrice.

\paragraph{Deux désaccords sur le sens des arcs, déclarés.} Il conteste que la gouvernance soit une
source pure : un incident significatif peut conduire à changer les responsabilités, à renforcer
les comités de suivi ou à augmenter les budgets, si bien que la gouvernance est aussi une
conséquence des défaillances. Il conteste de même que la gestion des incidents soit un puits : un
incident mal géré révèle des tests insuffisants et conduit à les renforcer, un incident impliquant
un prestataire fait revoir les contrats et la surveillance, et une escalade défaillante dégrade la
circulation de l'information. Deux lectures de ces exemples doivent être tenues ensemble. La
plupart décrivent une rétroaction corrective, l'incident conduisant à renforcer un domaine :
elle relève de l'évolution des états de conformité modélisée au chapitre~\ref{chap:multietats},
non de la propagation d'une perte que porte la matrice. Un seul décrit une dégradation en retour,
de la gestion des incidents vers le partage d'information, sens qu'aucun rapport du corpus ne
code. Ce désaccord ne déplace pas le capital publié : l'ensemble admissible du
chapitre~\ref{chap:partielle} ne fixe aucune direction et contient aussi bien les matrices à
dépendance réciproque que celles où l'arc est inversé. Il affaiblit en revanche la structure que
le corpus suggère, déjà déclarée corroborée et non calibrée, et il est enregistré comme tel.

\paragraph{La faisabilité, sous conditions.} Il juge les trois indicateurs pertinents et
actionnables mais insuffisants à eux seuls, et y ajouterait les incidents récurrents, les
vulnérabilités critiques non corrigées, les résultats des tests et la disponibilité des services
critiques. Il juge le registre d'incidents réaliste, et la distinction entre date de survenue et
date de déclaration particulièrement utile, sous réserve de règles de qualification précises :
le classement d'un incident par domaine peut être subjectif, et reconnaître une cause commune à
plusieurs victimes suppose de disposer de l'information qui l'établit. Sur le prestataire
partagé, il confirme le choc simultané en rappelant que le niveau de dépendance résulte lui-même
de choix de gouvernance.

\paragraph{Ce qui manque selon lui.} La gestion des changements et celle des vulnérabilités,
qu'il désigne comme des déclencheurs fréquents d'une chaîne traversant plusieurs domaines sans
qu'aucun ait été défaillant au départ. Le modèle fait naître un incident sur un pilier ; un
déclencheur transversal de ce type n'y a pas de place propre, et c'est une limite de périmètre à
inscrire parmi les prolongements.

\paragraph{Statut.} Deux réponses ne font pas un panel. Cette relecture n'entre donc dans aucune
table de robustesse et ne modifie aucun niveau publié. Elle est citée pour ce qu'elle est, une
relecture de praticiens engagée et partiellement rendue, et les répondants ne sont pas nommés.

\section{Le coût de la décision et sa réversibilité}

Il faut nommer le coût. L'élicitation avait une fonction précise et une seule :
\textbf{départager P1 et P4} en tête du classement des piliers sources. Sans elle, les deux
restent pratiquement à égalité sur l'ensemble admissible, P1 arrivant premier dans
$43{,}6\,\%$ des configurations et P4 dans $37{,}7\,\%$, pour un écart de regret maximal de
$1{,}5\,\%$ (\S\ref{sec:res-prio}). Le mémoire conclut donc que P1 et P4 devancent nettement les
trois autres piliers, et il s'arrête là où la donnée s'arrête, au lieu de trancher entre eux par
un avis.
% SOURCES-SCRIPTS: 30 31

\begin{cle}
\textbf{Une décision, pas un renoncement.} Le protocole est écrit, le questionnaire est prêt, le
gabarit de réponses et la chaîne d'agrégation sont testés. Ce qui manque n'est pas le
dispositif mais un panel de taille suffisante pour que la pondération de Cooke ait un sens.
La décision est donc \emph{réversible sans coût de développement} : si un accès à un panel
d'une dizaine de répondants devient possible, l'élicitation reprend là où elle a été laissée,
et son apport attendu est connu d'avance, à savoir le départage de P1 et P4 et un resserrement
de la bande d'identification. C'est le seul point du mémoire où une donnée supplémentaire
changerait une conclusion plutôt qu'un niveau.
\end{cle}
